\scanpage{9}
\mspage{8}

\deleted{\uline{Définition} Un \deleted{\unsure{groupe est dit}} \inserted{préschéma en} groupe \inserted{\deleted{\unsure{de type fini}}} $X$ sur $S$ est dit $\pi_0$-rigide}

\uline{Définition} D'une condition sur un préschéma en groupe $G$ de \deleted{type fini}\inserted{\deleted{plat}} sur $S$: \deleted{\unsure{local noeth et}} pour tt schém. local $T$ sur $S$, et tout sous-groupe ouvert \deleted{et fermé} $U$ de $G_T$ qui majore la fibre spéciale, on a $U=G_T$.

\margin{\uline{Condition R}}

Lorsque \deleted{\unsure{$G$ sur}} \deleted{\unsure{est}} $G$ est de type fini sur $S$ loc. noeth, il suffit de tester aux $T$ spectr. d'un anneau de valuation discrète [complet, corps rés. alg. clos].

On trouve donc

\uline{Corollaire} Soit $\unsure{X|G}$ de \deleted{\unsure{type plat ou présentation}}\inserted{type} fini sur $S$\inserted{loc noeth}, \deleted{$\unsure{H}$} $Y$ un sous-groupe fermé de, \deleted{présentation finie}. Supposons que $X$ satisfasse R. Alors $\prod_{X/S}Y/X$ est représentable par un sous-préschéma fermé de $S$.

Exemples

\deleted{Soit}

\deleted{\uline{Corollaire} Soit $G$ schém. en groupes noeth sur $S$ loc. noeth, $H$ un sous-groupe fermé de type fini sur $S$, satisfaisant la condition R. Alors $\underline{\operatorname{Cent.}}_G(H)$ \inserted{(si $G$ \unsure{séparé sur $S$})} et $\underline{\operatorname{Norm}}_G(H)$ sont représentables par un sous-préschéma fermé de $G$. De $=$ $\operatorname{Transp}_G(H,K)$}
